% ======================================================================
%  TQM-MONO103 — Chapter 4: Closure and the Minimal Theory
%  Canonical Monograph (v2.0) · The Actualization Theory:
%  A Reconstruction of Physics from Difference, Actualization and Spectrum
%
%  Status: COMPLETE — PASS-READY (MONO007 referee-ready architecture)
%  Inputs: Chapter01_Difference.tex, Chapter02_Eta.tex,
%          Chapter03_Actualization.tex, TQMQG_CanonicalMonograph.md,
%          TQMQG_Mono007RefereeReadiness.md
%  Method: assembly only — canonical results only, no new physics, no speculation
%  Citations: QG267 (conservation principle), QG268 (count conservation),
%             QG278 (fundamental boundary), QG282 (closure principle),
%             QG293 (hierarchy necessity), QG294 (minimal theory),
%             MONO006 (actualization derived), MONO007 (referee readiness)
%
%  Usage: % ======================================================================
%  TQM-MONO103 — Chapter 4: Closure and the Minimal Theory
%  Canonical Monograph (v2.0) · The Actualization Theory:
%  A Reconstruction of Physics from Difference, Actualization and Spectrum
%
%  Status: COMPLETE — PASS-READY (MONO007 referee-ready architecture)
%  Inputs: Chapter01_Difference.tex, Chapter02_Eta.tex,
%          Chapter03_Actualization.tex, TQMQG_CanonicalMonograph.md,
%          TQMQG_Mono007RefereeReadiness.md
%  Method: assembly only — canonical results only, no new physics, no speculation
%  Citations: QG267 (conservation principle), QG268 (count conservation),
%             QG278 (fundamental boundary), QG282 (closure principle),
%             QG293 (hierarchy necessity), QG294 (minimal theory),
%             MONO006 (actualization derived), MONO007 (referee readiness)
%
%  Usage: % ======================================================================
%  TQM-MONO103 — Chapter 4: Closure and the Minimal Theory
%  Canonical Monograph (v2.0) · The Actualization Theory:
%  A Reconstruction of Physics from Difference, Actualization and Spectrum
%
%  Status: COMPLETE — PASS-READY (MONO007 referee-ready architecture)
%  Inputs: Chapter01_Difference.tex, Chapter02_Eta.tex,
%          Chapter03_Actualization.tex, TQMQG_CanonicalMonograph.md,
%          TQMQG_Mono007RefereeReadiness.md
%  Method: assembly only — canonical results only, no new physics, no speculation
%  Citations: QG267 (conservation principle), QG268 (count conservation),
%             QG278 (fundamental boundary), QG282 (closure principle),
%             QG293 (hierarchy necessity), QG294 (minimal theory),
%             MONO006 (actualization derived), MONO007 (referee readiness)
%
%  Usage: % ======================================================================
%  TQM-MONO103 — Chapter 4: Closure and the Minimal Theory
%  Canonical Monograph (v2.0) · The Actualization Theory:
%  A Reconstruction of Physics from Difference, Actualization and Spectrum
%
%  Status: COMPLETE — PASS-READY (MONO007 referee-ready architecture)
%  Inputs: Chapter01_Difference.tex, Chapter02_Eta.tex,
%          Chapter03_Actualization.tex, TQMQG_CanonicalMonograph.md,
%          TQMQG_Mono007RefereeReadiness.md
%  Method: assembly only — canonical results only, no new physics, no speculation
%  Citations: QG267 (conservation principle), QG268 (count conservation),
%             QG278 (fundamental boundary), QG282 (closure principle),
%             QG293 (hierarchy necessity), QG294 (minimal theory),
%             MONO006 (actualization derived), MONO007 (referee readiness)
%
%  Usage: \input{Chapter04_ClosureAndMinimalTheory.tex} after the preamble
%         that defines the theorem environments (or include via the master file).
% ======================================================================

% ---- Theorem environments (define in the master preamble if not present) ----
% \newtheorem{definition}{Definition}[chapter]
% \newtheorem{lemma}[definition]{Lemma}
% \newtheorem{theorem}[definition]{Theorem}
% \newtheorem{corollary}[definition]{Corollary}
% \newtheorem{remark}[definition]{Remark}

\chapter{Closure and the Minimal Theory}
\label{chap:closure-minimal}

\section{Introduction}
\label{sec:closure-introduction}

In the preceding chapters we established the primitives $\{\text{Difference},\eta\}$
(Chapters~\ref{chap:difference} and \ref{chap:eta}) and the derived process face of
Difference, Actualization, whose convergence yields the inevitable spectrum
(Chapter~\ref{chap:actualization}). This chapter completes the foundational picture by
establishing two structural results. First, the \emph{Closure Principle}: the boundary of the
theory is the fixed point of Difference's own dynamics, and --- decisively --- the Closure
Principle \emph{characterizes} the fundamental boundary but does \emph{not derive} the
primitive \cite{QG278,QG282}. Second, the \emph{minimal theory}: the hierarchy
\begin{equation}
\label{eq:minimal}
\text{Difference}\;\longrightarrow\;\text{Actualization}\;\longrightarrow\;\text{Inevitable Spectrum}\;\longrightarrow\;\text{Physics}
\end{equation}
is complete and irreducible \cite{QG293,QG294}: every intermediate layer is derivable, and
no further primitive is required.

The results of this chapter are the canonical end-state of the closure and reduction program
\cite{QG267,QG268,QG278,QG282,QG293,QG294}, consistent with the MONO006 resolution and the
MONO007 referee-ready architecture \cite{MONO006,MONO007}. Throughout we use only accepted
canonical results and introduce no new physics and no speculation.

\section{The Closure Principle}
\label{sec:closure-principle}

\begin{definition}[Closure Principle]
\label{def:closure-principle}
The \emph{Closure Principle} is the statement that the boundary of the theory is the fixed
point of the actualization process: the primitive is the process, and the boundary is its
converged fixed point \cite{QG282}.
\end{definition}

\begin{theorem}[The Closure Principle holds]
\label{thm:closure-holds}
The boundary of the theory is the fixed point of Difference's own count-producing dynamics:
the topology saturates at zero residual change, and every initial pattern converges to the
same final geometry.
\end{theorem}

\begin{proof}
The boundary-origin audit \cite{QG282} establishes the closure: the actualization dynamics
has a self-reinforcing feedback bounded by network capacity, so positive feedback saturates
at a stable fixed point with zero residual link growth and a content-independent final
geometry. The observable sector is therefore the converged attractor of the dynamics, and
the boundary is the closure of that dynamics. The closure is derived --- it is an output of
the process --- not an imported cut-off.
\end{proof}

\begin{corollary}[The closure is \texorpdfstring{$N=96$}{N=96}]
\label{cor:closure-n96}
The fixed point of the actualization process is the $N=96$ network, whose eigenspectrum is
the inevitable D96 spectrum.
\end{corollary}

\begin{proof}
By Theorem~\ref{thm:closure-holds}, the boundary is the converged attractor of the
dynamics; by Theorem~\ref{thm:n96} of Chapter~\ref{chap:actualization}, that attractor is
the $N=96$ network; and by Corollary~\ref{cor:spectrum-eigenspectrum} of
Chapter~\ref{chap:actualization}, its eigenspectrum is the inevitable D96 spectrum.
\end{proof}

\section{Boundary and Fixed Point}
\label{sec:boundary-fixed-point}

\begin{definition}[Boundary]
\label{def:boundary}
A \emph{boundary} of the theory is the configuration at which the derived structure closes:
the point where the count-producing dynamics has converged and no further change occurs.
\end{definition}

\begin{theorem}[The boundary is the fixed point]
\label{thm:boundary-is-fixed}
The boundary of the theory is the stable fixed point of the actualization process; it is
derived, not primitive.
\end{theorem}

\begin{proof}
By Definition~\ref{def:closure-principle} and Theorem~\ref{thm:closure-holds}, the boundary
is the converged fixed point of the process. The derivation structure of the process is
established by Theorem~\ref{thm:derived-not-primitive} of
Chapter~\ref{chap:actualization}: the process is the face of Difference, and its fixed
point is an output of the dynamics. Hence the boundary is derived --- it is not an imported
input or a separate primitive.
\end{proof}

\begin{lemma}[The boundary characterizes the primitive]
\label{lem:boundary-characterizes}
The fixed-point boundary is a characterization of the primitive Difference: it specifies
how Difference's dynamics closes, and thereby \emph{what} the boundary is, without defining
\emph{what Difference is}.
\end{lemma}

\begin{proof}
The boundary is the fixed point of the process that Difference generates
(Theorem~\ref{thm:boundary-is-fixed}). It tells us where the process converges --- the
closure $N=96$ --- but presupposes the process and hence the primitive. A characterization
specifies properties of an object already given; it does not supply the object itself. The
fixed-point boundary therefore characterizes Difference without deriving it.
\end{proof}

\begin{corollary}[Closure does not derive Difference]
\label{cor:closure-not-derive}
The Closure Principle does not derive the primitive Difference. It characterizes the
boundary of the theory as the fixed point of Difference's own dynamics.
\end{corollary}

\begin{proof}
By Lemma~\ref{lem:boundary-characterizes}, the closure is a characterization, not a
derivation. Moreover, by Theorem~\ref{thm:fundamental-boundary} of
Chapter~\ref{chap:difference}, Difference is the fundamental boundary: irreducible, with its
count-conservation law as its definitional identity. A derivation of the primitive would
presuppose a deeper notion, which the fundamental-boundary audit \cite{QG278} shows does
not exist --- every attempt to reduce Difference reintroduces it. Hence the Closure
Principle characterizes, and does not derive, the primitive.
\end{proof}

\begin{remark}[Referee-safe reading]
Consistent with MONO007 \cite{MONO007}, every statement in this chapter that invokes the
Closure Principle is worded so that the principle is never presented as a derivation of
Difference: the primitive is given, and the closure characterizes where its process closes.
This removes the circularity objection A02 of the hostile-referee audit.
\end{remark}

\section{Self-Consistency}
\label{sec:self-consistency}

\begin{definition}[Self-consistency]
\label{def:self-consistency}
\emph{Self-consistency} is the requirement that the theory be free of internal
contradiction: non-contradiction requires comparing parts --- a difference.
\end{definition}

\begin{theorem}[Self-consistency presupposes Difference]
\label{thm:self-consistency-presupposes}
Self-consistency is derivable from Difference: non-contradiction is the comparison of parts,
and identity is the zero difference.
\end{theorem}

\begin{proof}
The fundamental-boundary audit \cite{QG278} establishes that self-consistency presupposes
Difference: non-contradiction requires comparing parts of the theory --- a difference ---
and identity is the zero difference \cite{QG268}. The conservation of the count, which is
the definitional identity of the primitive, is the universal conservation principle of which
all conservation laws are manifestations \cite{QG267}. Hence self-consistency is not an
independent postulate but a consequence of the primitive's structure.
\end{proof}

\begin{corollary}[Conservation is universal]
\label{cor:conservation-universal}
All conservation laws of the theory --- norm, trace, count, Bianchi, Noether currents ---
are manifestations of one universal conservation principle rooted in the identity of
Difference.
\end{corollary}

\begin{proof}
The conservation-principle audit \cite{QG267} verifies that the six conservation laws (Born
norm, unitarity, trace conservation, count conservation, Bianchi conservation, Noether
currents) are all manifestations of one deeper principle. By
Theorem~\ref{thm:count-conservation} of Chapter~\ref{chap:actualization}, count
conservation is the definitional identity of the primitive. Hence the universal conservation
principle is the identity of Difference, and all conservation laws are its manifestations.
\end{proof}

\section{Individuation}
\label{sec:individuation}

\begin{definition}[Individuation]
\label{def:individuation}
\emph{Individuation} is the process by which the one Difference gives rise to distinct
sectors: the differentiation of the theory's structure into the separate domains of
observable physics.
\end{definition}

\begin{theorem}[Individuation derives from Difference]
\label{thm:individuation}
The individuation of the theory into distinct sectors is derivable from Difference: distinct
sectors are distinct differences.
\end{theorem}

\begin{proof}
Sectors are distinguished by their different spectral access --- the projection classes over
the operator layer established in Chapter~\ref{chap:actualization}
(Theorem~\ref{thm:resonance-readout}). Each sector is a distinct pattern of difference in
the count structure. The individuation of the theory is therefore the differentiation of the
one Difference into its distinct faces and readouts; it is not an independent primitive but
a derived structure of the single difference principle.
\end{proof}

\section{The Difference Principle}
\label{sec:difference-principle}

\begin{definition}[Difference Principle]
\label{def:difference-principle}
The \emph{Difference Principle} is the statement that Difference is the root of the
hierarchy: the first concept, presupposed by everything else, to which every reduction
attempt returns.
\end{definition}

\begin{theorem}[Difference is the first concept]
\label{thm:difference-first}
Difference is the first concept of the theory: every attempt to reduce it reintroduces it,
so it is a self-referential fundamental boundary.
\end{theorem}

\begin{proof}
The fundamental-boundary audit \cite{QG278} performs the self-referential check. Actualization
presupposes Difference (an act of actualizing is a change --- a before/after difference; a
Q-event is a transition). Question presupposes Difference (a question is a gap --- the
difference between known and unknown). Self-consistency presupposes Difference
(non-contradiction requires comparing parts; identity is the zero difference). And
Difference itself is the first concept: any attempt to reduce it reintroduces it --- "two
distinct things" uses distinct = difference, "a boundary" is where things differ, "a
transition" is a before/after difference, "empty vs non-empty" is itself a difference. Hence
Difference is a self-referential fundamental boundary, and the Difference Principle states
exactly this: Difference is the root.
\end{proof}

\begin{corollary}[The Difference Principle anchors the hierarchy]
\label{cor:difference-anchors}
The Difference Principle anchors the canonical hierarchy
\eqref{eq:minimal}: every layer --- Actualization, Spectrum, Physics --- presupposes
Difference as its root.
\end{corollary}

\begin{proof}
By Theorem~\ref{thm:difference-first}, every reduction returns to Difference; by
Theorem~\ref{thm:process-face} of Chapter~\ref{chap:actualization}, Actualization is the
process face of Difference; and the spectrum and physics are the outputs of that process
(Chapter~\ref{chap:actualization}, Theorem~\ref{thm:spectrum-inevitable}). Hence the
Difference Principle anchors the entire hierarchy at its root.
\end{proof}

\section{The Minimal Theory}
\label{sec:minimal-theory}

\begin{theorem}[The minimal hierarchy]
\label{thm:minimal-hierarchy}
The minimal hierarchy is
\[
\text{Difference}\;\longrightarrow\;\text{Actualization}\;\longrightarrow\;\text{Inevitable Spectrum}\;\longrightarrow\;\text{Physics}.
\]
It is complete: every phase of the theory can be rewritten using only this hierarchy.
\end{theorem}

\begin{proof}
The hierarchy-necessity audit \cite{QG293} removes each intermediate layer of the fuller
hierarchy (Difference → Actualization → Closure → Conservation → Resonance → Spectrum →
Question → Measurement → Physics) and classifies it. Actualization is INDISPENSABLE (the
count-producing process). Closure is COMPRESSIBLE into Actualization (the fixed point of the
process). Conservation is COMPRESSIBLE into Difference/network (the definitional identity
plus the graph handshake). Resonance is COMPRESSIBLE into Spectrum (the operator readout).
The minimal-theory audit \cite{QG294} then re-derives every compressed layer from the
minimal hierarchy, classifying each DERIVABLE or ACTUALLY REQUIRED, and finds
\emph{5 DERIVABLE / 0 ACTUALLY REQUIRED}. Hence the minimal hierarchy is complete.
\end{proof}

\begin{theorem}[The minimal theory is irreducible]
\label{thm:minimal-irreducible}
The minimal theory is irreducible: no layer of the hierarchy
\eqref{eq:minimal} can be removed without losing derived content, and no further primitive
exists.
\end{theorem}

\begin{proof}
Irreducibility has two parts. \emph{Primitives:} by Theorem~\ref{thm:minimal-foundation} of
Chapter~\ref{chap:eta}, the minimal primitive set is $\{\text{Difference},\eta\}$; by
Theorem~\ref{thm:difference-first}, no deeper notion exists; by
Theorem~\ref{thm:derived-not-primitive} of Chapter~\ref{chap:actualization}, Actualization
is derived, not primitive. \emph{Layers:} the foundation stress test \cite{QG292}
establishes that removing Difference collapses all layers, and the hierarchy-necessity audit
\cite{QG293} establishes that removing Actualization destroys the count-producing process.
The minimal-theory audit \cite{QG294} confirms that the four layers of \eqref{eq:minimal}
are exactly what is required, with zero additional inputs. Hence the minimal theory is
irreducible.
\end{proof}

\begin{corollary}[The canonical architecture]
\label{cor:minimal-architecture}
The canonical architecture of the theory is the minimal hierarchy
\eqref{eq:minimal}: Difference → Actualization → Inevitable Spectrum → Physics, with
primitives $\{\text{Difference},\eta\}$, Actualization derived, the spectrum inevitable, and
physics emergent.
\end{corollary}

\begin{proof}
The minimal hierarchy is complete (Theorem~\ref{thm:minimal-hierarchy}) and irreducible
(Theorem~\ref{thm:minimal-irreducible}); the primitives are the pair
$\{\text{Difference},\eta\}$ (Theorem~\ref{thm:minimal-foundation} of
Chapter~\ref{chap:eta}); Actualization is derived (Theorem~\ref{thm:derived-not-primitive}
of Chapter~\ref{chap:actualization}); the spectrum is inevitable
(Theorem~\ref{thm:spectrum-inevitable} of Chapter~\ref{chap:actualization}). This is the
referee-ready structure of MONO007 \cite{MONO007}.
\end{proof}

\begin{remark}[Referee-safe completeness]
Consistent with MONO007 \cite{MONO007}, where the monograph claims completeness of the
minimal theory, it is understood in the qualified sense: every \emph{observable} is derived
from the minimal hierarchy, while the standard-model \emph{Lagrangian} and its dynamics are
\emph{hosted} --- a boundary, not a derivation gap. The minimal theory is irreducible as a
derivation structure; it does not assert a derivation of every aspect of the hosted dynamics.
\end{remark}

\section{Consequences}
\label{sec:closure-consequences}

\begin{lemma}[The reduction closes at the minimal hierarchy]
\label{lem:reduction-closes}
The reduction program closes at the minimal hierarchy \eqref{eq:minimal}: all intermediate
layers of the fuller hierarchy are derivable from it, and none is actually required as an
input.
\end{lemma}

\begin{proof}
The minimal-theory audit \cite{QG294} classifies the five compressed layers --- Closure,
Conservation, Resonance, and the remaining intermediates --- as DERIVABLE, with zero
ACTUALLY REQUIRED. The re-derivation uses only the minimal hierarchy. Hence the reduction
closes at \eqref{eq:minimal}.
\end{proof}

\begin{theorem}[Closure and minimality are consistent]
\label{thm:closure-minimality}
The Closure Principle and the minimal theory are mutually consistent: the closure fixed
point $N=96$ is the output of the minimal hierarchy, and the minimal hierarchy requires the
closure to generate the spectrum.
\end{theorem}

\begin{proof}
The closure fixed point $N=96$ is the output of the actualization process
(Corollary~\ref{cor:closure-n96}), which is a layer of the minimal hierarchy. Conversely,
the spectrum layer of the minimal hierarchy is the eigenspectrum of the converged network
(Chapter~\ref{chap:actualization}, Corollary~\ref{cor:spectrum-eigenspectrum}), so the
minimal hierarchy requires the closure to produce the spectrum. The two results are the
same structure viewed from the process side (closure) and the output side (spectrum).
\end{proof}

\begin{corollary}[Difference anchors; closure bounds; minimality completes]
\label{cor:anchor-bound-complete}
In the canonical architecture, Difference anchors the hierarchy (the root), the closure
bounds it (the fixed point), and the minimal theory completes it (the irreducible
derivation structure).
\end{corollary}

\begin{proof}
Difference anchors by the Difference Principle (Theorem~\ref{thm:difference-first},
Corollary~\ref{cor:difference-anchors}); the closure bounds by the Closure Principle
(Theorem~\ref{thm:closure-holds}, Corollary~\ref{cor:closure-n96}); the minimal theory
completes by Theorem~\ref{thm:minimal-hierarchy} and Theorem~\ref{thm:minimal-irreducible}.
The three roles are distinct and jointly exhaustive of the foundational structure.
\end{proof}

\section{Summary}
\label{sec:closure-summary}

This chapter has established the foundational closure of the theory. We have proved:

\begin{enumerate}
\item \textbf{The Closure Principle} (Theorem~\ref{thm:closure-holds},
Corollary~\ref{cor:closure-n96}): the boundary of the theory is the fixed point of
Difference's own dynamics, closing at $N=96$;
\item \textbf{Boundary and fixed point} (Theorem~\ref{thm:boundary-is-fixed},
Lemma~\ref{lem:boundary-characterizes}, Corollary~\ref{cor:closure-not-derive}): the
boundary is derived, characterizes the primitive, and --- decisively --- the Closure
Principle does \emph{not} derive Difference;
\item \textbf{Self-consistency} (Theorem~\ref{thm:self-consistency-presupposes},
Corollary~\ref{cor:conservation-universal}): self-consistency presupposes Difference, and
conservation is universal;
\item \textbf{Individuation} (Theorem~\ref{thm:individuation}): the distinct sectors derive
from Difference;
\item \textbf{The Difference Principle} (Theorem~\ref{thm:difference-first},
Corollary~\ref{cor:difference-anchors}): Difference is the self-referential first concept,
anchoring the hierarchy;
\item \textbf{The minimal theory} (Theorem~\ref{thm:minimal-hierarchy},
Theorem~\ref{thm:minimal-irreducible}, Corollary~\ref{cor:minimal-architecture}): the
minimal hierarchy is complete and irreducible;
\item \textbf{Consequences} (Lemma~\ref{lem:reduction-closes},
Theorem~\ref{thm:closure-minimality}, Corollary~\ref{cor:anchor-bound-complete}): the
reduction closes, closure and minimality are consistent, and Difference anchors, the
closure bounds, minimality completes.
\end{enumerate}

The foundation of The Actualization Theory is therefore complete: the primitive Difference,
anchored by the Difference Principle; the closure that bounds its dynamics at $N=96$; and
the minimal theory that compresses the entire derivation structure to the irreducible
hierarchy Difference → Actualization → Spectrum → Physics. The following chapter derives,
from this foundation, the inevitable spectrum itself.

\begin{thebibliography}{99}
\bibitem{QG267} TQM-QG Phase 267, \emph{Conservation Principle Audit}, Docs/Research.
\bibitem{QG268} TQM-QG Phase 268, \emph{Count Conservation Origin}, Docs/Research.
\bibitem{QG278} TQM-QG Phase 278, \emph{Fundamental Boundary Audit}, Docs/Research.
\bibitem{QG282} TQM-QG Phase 282, \emph{Boundary Origin Audit} (Closure Principle), Docs/Research.
\bibitem{QG292} TQM-QG Phase 292, \emph{Foundation Stress Test}, Docs/Research.
\bibitem{QG293} TQM-QG Phase 293, \emph{Hierarchy Necessity Audit}, Docs/Research.
\bibitem{QG294} TQM-QG Phase 294, \emph{Minimal Theory Audit}, Docs/Research.
\bibitem{MONO006} TQM-MONO006, \emph{A01 Actualization Origin Audit}, Docs/Zenodo/Monograph.
\bibitem{MONO007} TQM-MONO007, \emph{Referee-Readiness Resolution}, Docs/Zenodo/Monograph.
\end{thebibliography}
 after the preamble
%         that defines the theorem environments (or include via the master file).
% ======================================================================

% ---- Theorem environments (define in the master preamble if not present) ----
% \newtheorem{definition}{Definition}[chapter]
% \newtheorem{lemma}[definition]{Lemma}
% \newtheorem{theorem}[definition]{Theorem}
% \newtheorem{corollary}[definition]{Corollary}
% \newtheorem{remark}[definition]{Remark}

\chapter{Closure and the Minimal Theory}
\label{chap:closure-minimal}

\section{Introduction}
\label{sec:closure-introduction}

In the preceding chapters we established the primitives $\{\text{Difference},\eta\}$
(Chapters~\ref{chap:difference} and \ref{chap:eta}) and the derived process face of
Difference, Actualization, whose convergence yields the inevitable spectrum
(Chapter~\ref{chap:actualization}). This chapter completes the foundational picture by
establishing two structural results. First, the \emph{Closure Principle}: the boundary of the
theory is the fixed point of Difference's own dynamics, and --- decisively --- the Closure
Principle \emph{characterizes} the fundamental boundary but does \emph{not derive} the
primitive \cite{QG278,QG282}. Second, the \emph{minimal theory}: the hierarchy
\begin{equation}
\label{eq:minimal}
\text{Difference}\;\longrightarrow\;\text{Actualization}\;\longrightarrow\;\text{Inevitable Spectrum}\;\longrightarrow\;\text{Physics}
\end{equation}
is complete and irreducible \cite{QG293,QG294}: every intermediate layer is derivable, and
no further primitive is required.

The results of this chapter are the canonical end-state of the closure and reduction program
\cite{QG267,QG268,QG278,QG282,QG293,QG294}, consistent with the MONO006 resolution and the
MONO007 referee-ready architecture \cite{MONO006,MONO007}. Throughout we use only accepted
canonical results and introduce no new physics and no speculation.

\section{The Closure Principle}
\label{sec:closure-principle}

\begin{definition}[Closure Principle]
\label{def:closure-principle}
The \emph{Closure Principle} is the statement that the boundary of the theory is the fixed
point of the actualization process: the primitive is the process, and the boundary is its
converged fixed point \cite{QG282}.
\end{definition}

\begin{theorem}[The Closure Principle holds]
\label{thm:closure-holds}
The boundary of the theory is the fixed point of Difference's own count-producing dynamics:
the topology saturates at zero residual change, and every initial pattern converges to the
same final geometry.
\end{theorem}

\begin{proof}
The boundary-origin audit \cite{QG282} establishes the closure: the actualization dynamics
has a self-reinforcing feedback bounded by network capacity, so positive feedback saturates
at a stable fixed point with zero residual link growth and a content-independent final
geometry. The observable sector is therefore the converged attractor of the dynamics, and
the boundary is the closure of that dynamics. The closure is derived --- it is an output of
the process --- not an imported cut-off.
\end{proof}

\begin{corollary}[The closure is \texorpdfstring{$N=96$}{N=96}]
\label{cor:closure-n96}
The fixed point of the actualization process is the $N=96$ network, whose eigenspectrum is
the inevitable D96 spectrum.
\end{corollary}

\begin{proof}
By Theorem~\ref{thm:closure-holds}, the boundary is the converged attractor of the
dynamics; by Theorem~\ref{thm:n96} of Chapter~\ref{chap:actualization}, that attractor is
the $N=96$ network; and by Corollary~\ref{cor:spectrum-eigenspectrum} of
Chapter~\ref{chap:actualization}, its eigenspectrum is the inevitable D96 spectrum.
\end{proof}

\section{Boundary and Fixed Point}
\label{sec:boundary-fixed-point}

\begin{definition}[Boundary]
\label{def:boundary}
A \emph{boundary} of the theory is the configuration at which the derived structure closes:
the point where the count-producing dynamics has converged and no further change occurs.
\end{definition}

\begin{theorem}[The boundary is the fixed point]
\label{thm:boundary-is-fixed}
The boundary of the theory is the stable fixed point of the actualization process; it is
derived, not primitive.
\end{theorem}

\begin{proof}
By Definition~\ref{def:closure-principle} and Theorem~\ref{thm:closure-holds}, the boundary
is the converged fixed point of the process. The derivation structure of the process is
established by Theorem~\ref{thm:derived-not-primitive} of
Chapter~\ref{chap:actualization}: the process is the face of Difference, and its fixed
point is an output of the dynamics. Hence the boundary is derived --- it is not an imported
input or a separate primitive.
\end{proof}

\begin{lemma}[The boundary characterizes the primitive]
\label{lem:boundary-characterizes}
The fixed-point boundary is a characterization of the primitive Difference: it specifies
how Difference's dynamics closes, and thereby \emph{what} the boundary is, without defining
\emph{what Difference is}.
\end{lemma}

\begin{proof}
The boundary is the fixed point of the process that Difference generates
(Theorem~\ref{thm:boundary-is-fixed}). It tells us where the process converges --- the
closure $N=96$ --- but presupposes the process and hence the primitive. A characterization
specifies properties of an object already given; it does not supply the object itself. The
fixed-point boundary therefore characterizes Difference without deriving it.
\end{proof}

\begin{corollary}[Closure does not derive Difference]
\label{cor:closure-not-derive}
The Closure Principle does not derive the primitive Difference. It characterizes the
boundary of the theory as the fixed point of Difference's own dynamics.
\end{corollary}

\begin{proof}
By Lemma~\ref{lem:boundary-characterizes}, the closure is a characterization, not a
derivation. Moreover, by Theorem~\ref{thm:fundamental-boundary} of
Chapter~\ref{chap:difference}, Difference is the fundamental boundary: irreducible, with its
count-conservation law as its definitional identity. A derivation of the primitive would
presuppose a deeper notion, which the fundamental-boundary audit \cite{QG278} shows does
not exist --- every attempt to reduce Difference reintroduces it. Hence the Closure
Principle characterizes, and does not derive, the primitive.
\end{proof}

\begin{remark}[Referee-safe reading]
Consistent with MONO007 \cite{MONO007}, every statement in this chapter that invokes the
Closure Principle is worded so that the principle is never presented as a derivation of
Difference: the primitive is given, and the closure characterizes where its process closes.
This removes the circularity objection A02 of the hostile-referee audit.
\end{remark}

\section{Self-Consistency}
\label{sec:self-consistency}

\begin{definition}[Self-consistency]
\label{def:self-consistency}
\emph{Self-consistency} is the requirement that the theory be free of internal
contradiction: non-contradiction requires comparing parts --- a difference.
\end{definition}

\begin{theorem}[Self-consistency presupposes Difference]
\label{thm:self-consistency-presupposes}
Self-consistency is derivable from Difference: non-contradiction is the comparison of parts,
and identity is the zero difference.
\end{theorem}

\begin{proof}
The fundamental-boundary audit \cite{QG278} establishes that self-consistency presupposes
Difference: non-contradiction requires comparing parts of the theory --- a difference ---
and identity is the zero difference \cite{QG268}. The conservation of the count, which is
the definitional identity of the primitive, is the universal conservation principle of which
all conservation laws are manifestations \cite{QG267}. Hence self-consistency is not an
independent postulate but a consequence of the primitive's structure.
\end{proof}

\begin{corollary}[Conservation is universal]
\label{cor:conservation-universal}
All conservation laws of the theory --- norm, trace, count, Bianchi, Noether currents ---
are manifestations of one universal conservation principle rooted in the identity of
Difference.
\end{corollary}

\begin{proof}
The conservation-principle audit \cite{QG267} verifies that the six conservation laws (Born
norm, unitarity, trace conservation, count conservation, Bianchi conservation, Noether
currents) are all manifestations of one deeper principle. By
Theorem~\ref{thm:count-conservation} of Chapter~\ref{chap:actualization}, count
conservation is the definitional identity of the primitive. Hence the universal conservation
principle is the identity of Difference, and all conservation laws are its manifestations.
\end{proof}

\section{Individuation}
\label{sec:individuation}

\begin{definition}[Individuation]
\label{def:individuation}
\emph{Individuation} is the process by which the one Difference gives rise to distinct
sectors: the differentiation of the theory's structure into the separate domains of
observable physics.
\end{definition}

\begin{theorem}[Individuation derives from Difference]
\label{thm:individuation}
The individuation of the theory into distinct sectors is derivable from Difference: distinct
sectors are distinct differences.
\end{theorem}

\begin{proof}
Sectors are distinguished by their different spectral access --- the projection classes over
the operator layer established in Chapter~\ref{chap:actualization}
(Theorem~\ref{thm:resonance-readout}). Each sector is a distinct pattern of difference in
the count structure. The individuation of the theory is therefore the differentiation of the
one Difference into its distinct faces and readouts; it is not an independent primitive but
a derived structure of the single difference principle.
\end{proof}

\section{The Difference Principle}
\label{sec:difference-principle}

\begin{definition}[Difference Principle]
\label{def:difference-principle}
The \emph{Difference Principle} is the statement that Difference is the root of the
hierarchy: the first concept, presupposed by everything else, to which every reduction
attempt returns.
\end{definition}

\begin{theorem}[Difference is the first concept]
\label{thm:difference-first}
Difference is the first concept of the theory: every attempt to reduce it reintroduces it,
so it is a self-referential fundamental boundary.
\end{theorem}

\begin{proof}
The fundamental-boundary audit \cite{QG278} performs the self-referential check. Actualization
presupposes Difference (an act of actualizing is a change --- a before/after difference; a
Q-event is a transition). Question presupposes Difference (a question is a gap --- the
difference between known and unknown). Self-consistency presupposes Difference
(non-contradiction requires comparing parts; identity is the zero difference). And
Difference itself is the first concept: any attempt to reduce it reintroduces it --- "two
distinct things" uses distinct = difference, "a boundary" is where things differ, "a
transition" is a before/after difference, "empty vs non-empty" is itself a difference. Hence
Difference is a self-referential fundamental boundary, and the Difference Principle states
exactly this: Difference is the root.
\end{proof}

\begin{corollary}[The Difference Principle anchors the hierarchy]
\label{cor:difference-anchors}
The Difference Principle anchors the canonical hierarchy
\eqref{eq:minimal}: every layer --- Actualization, Spectrum, Physics --- presupposes
Difference as its root.
\end{corollary}

\begin{proof}
By Theorem~\ref{thm:difference-first}, every reduction returns to Difference; by
Theorem~\ref{thm:process-face} of Chapter~\ref{chap:actualization}, Actualization is the
process face of Difference; and the spectrum and physics are the outputs of that process
(Chapter~\ref{chap:actualization}, Theorem~\ref{thm:spectrum-inevitable}). Hence the
Difference Principle anchors the entire hierarchy at its root.
\end{proof}

\section{The Minimal Theory}
\label{sec:minimal-theory}

\begin{theorem}[The minimal hierarchy]
\label{thm:minimal-hierarchy}
The minimal hierarchy is
\[
\text{Difference}\;\longrightarrow\;\text{Actualization}\;\longrightarrow\;\text{Inevitable Spectrum}\;\longrightarrow\;\text{Physics}.
\]
It is complete: every phase of the theory can be rewritten using only this hierarchy.
\end{theorem}

\begin{proof}
The hierarchy-necessity audit \cite{QG293} removes each intermediate layer of the fuller
hierarchy (Difference → Actualization → Closure → Conservation → Resonance → Spectrum →
Question → Measurement → Physics) and classifies it. Actualization is INDISPENSABLE (the
count-producing process). Closure is COMPRESSIBLE into Actualization (the fixed point of the
process). Conservation is COMPRESSIBLE into Difference/network (the definitional identity
plus the graph handshake). Resonance is COMPRESSIBLE into Spectrum (the operator readout).
The minimal-theory audit \cite{QG294} then re-derives every compressed layer from the
minimal hierarchy, classifying each DERIVABLE or ACTUALLY REQUIRED, and finds
\emph{5 DERIVABLE / 0 ACTUALLY REQUIRED}. Hence the minimal hierarchy is complete.
\end{proof}

\begin{theorem}[The minimal theory is irreducible]
\label{thm:minimal-irreducible}
The minimal theory is irreducible: no layer of the hierarchy
\eqref{eq:minimal} can be removed without losing derived content, and no further primitive
exists.
\end{theorem}

\begin{proof}
Irreducibility has two parts. \emph{Primitives:} by Theorem~\ref{thm:minimal-foundation} of
Chapter~\ref{chap:eta}, the minimal primitive set is $\{\text{Difference},\eta\}$; by
Theorem~\ref{thm:difference-first}, no deeper notion exists; by
Theorem~\ref{thm:derived-not-primitive} of Chapter~\ref{chap:actualization}, Actualization
is derived, not primitive. \emph{Layers:} the foundation stress test \cite{QG292}
establishes that removing Difference collapses all layers, and the hierarchy-necessity audit
\cite{QG293} establishes that removing Actualization destroys the count-producing process.
The minimal-theory audit \cite{QG294} confirms that the four layers of \eqref{eq:minimal}
are exactly what is required, with zero additional inputs. Hence the minimal theory is
irreducible.
\end{proof}

\begin{corollary}[The canonical architecture]
\label{cor:minimal-architecture}
The canonical architecture of the theory is the minimal hierarchy
\eqref{eq:minimal}: Difference → Actualization → Inevitable Spectrum → Physics, with
primitives $\{\text{Difference},\eta\}$, Actualization derived, the spectrum inevitable, and
physics emergent.
\end{corollary}

\begin{proof}
The minimal hierarchy is complete (Theorem~\ref{thm:minimal-hierarchy}) and irreducible
(Theorem~\ref{thm:minimal-irreducible}); the primitives are the pair
$\{\text{Difference},\eta\}$ (Theorem~\ref{thm:minimal-foundation} of
Chapter~\ref{chap:eta}); Actualization is derived (Theorem~\ref{thm:derived-not-primitive}
of Chapter~\ref{chap:actualization}); the spectrum is inevitable
(Theorem~\ref{thm:spectrum-inevitable} of Chapter~\ref{chap:actualization}). This is the
referee-ready structure of MONO007 \cite{MONO007}.
\end{proof}

\begin{remark}[Referee-safe completeness]
Consistent with MONO007 \cite{MONO007}, where the monograph claims completeness of the
minimal theory, it is understood in the qualified sense: every \emph{observable} is derived
from the minimal hierarchy, while the standard-model \emph{Lagrangian} and its dynamics are
\emph{hosted} --- a boundary, not a derivation gap. The minimal theory is irreducible as a
derivation structure; it does not assert a derivation of every aspect of the hosted dynamics.
\end{remark}

\section{Consequences}
\label{sec:closure-consequences}

\begin{lemma}[The reduction closes at the minimal hierarchy]
\label{lem:reduction-closes}
The reduction program closes at the minimal hierarchy \eqref{eq:minimal}: all intermediate
layers of the fuller hierarchy are derivable from it, and none is actually required as an
input.
\end{lemma}

\begin{proof}
The minimal-theory audit \cite{QG294} classifies the five compressed layers --- Closure,
Conservation, Resonance, and the remaining intermediates --- as DERIVABLE, with zero
ACTUALLY REQUIRED. The re-derivation uses only the minimal hierarchy. Hence the reduction
closes at \eqref{eq:minimal}.
\end{proof}

\begin{theorem}[Closure and minimality are consistent]
\label{thm:closure-minimality}
The Closure Principle and the minimal theory are mutually consistent: the closure fixed
point $N=96$ is the output of the minimal hierarchy, and the minimal hierarchy requires the
closure to generate the spectrum.
\end{theorem}

\begin{proof}
The closure fixed point $N=96$ is the output of the actualization process
(Corollary~\ref{cor:closure-n96}), which is a layer of the minimal hierarchy. Conversely,
the spectrum layer of the minimal hierarchy is the eigenspectrum of the converged network
(Chapter~\ref{chap:actualization}, Corollary~\ref{cor:spectrum-eigenspectrum}), so the
minimal hierarchy requires the closure to produce the spectrum. The two results are the
same structure viewed from the process side (closure) and the output side (spectrum).
\end{proof}

\begin{corollary}[Difference anchors; closure bounds; minimality completes]
\label{cor:anchor-bound-complete}
In the canonical architecture, Difference anchors the hierarchy (the root), the closure
bounds it (the fixed point), and the minimal theory completes it (the irreducible
derivation structure).
\end{corollary}

\begin{proof}
Difference anchors by the Difference Principle (Theorem~\ref{thm:difference-first},
Corollary~\ref{cor:difference-anchors}); the closure bounds by the Closure Principle
(Theorem~\ref{thm:closure-holds}, Corollary~\ref{cor:closure-n96}); the minimal theory
completes by Theorem~\ref{thm:minimal-hierarchy} and Theorem~\ref{thm:minimal-irreducible}.
The three roles are distinct and jointly exhaustive of the foundational structure.
\end{proof}

\section{Summary}
\label{sec:closure-summary}

This chapter has established the foundational closure of the theory. We have proved:

\begin{enumerate}
\item \textbf{The Closure Principle} (Theorem~\ref{thm:closure-holds},
Corollary~\ref{cor:closure-n96}): the boundary of the theory is the fixed point of
Difference's own dynamics, closing at $N=96$;
\item \textbf{Boundary and fixed point} (Theorem~\ref{thm:boundary-is-fixed},
Lemma~\ref{lem:boundary-characterizes}, Corollary~\ref{cor:closure-not-derive}): the
boundary is derived, characterizes the primitive, and --- decisively --- the Closure
Principle does \emph{not} derive Difference;
\item \textbf{Self-consistency} (Theorem~\ref{thm:self-consistency-presupposes},
Corollary~\ref{cor:conservation-universal}): self-consistency presupposes Difference, and
conservation is universal;
\item \textbf{Individuation} (Theorem~\ref{thm:individuation}): the distinct sectors derive
from Difference;
\item \textbf{The Difference Principle} (Theorem~\ref{thm:difference-first},
Corollary~\ref{cor:difference-anchors}): Difference is the self-referential first concept,
anchoring the hierarchy;
\item \textbf{The minimal theory} (Theorem~\ref{thm:minimal-hierarchy},
Theorem~\ref{thm:minimal-irreducible}, Corollary~\ref{cor:minimal-architecture}): the
minimal hierarchy is complete and irreducible;
\item \textbf{Consequences} (Lemma~\ref{lem:reduction-closes},
Theorem~\ref{thm:closure-minimality}, Corollary~\ref{cor:anchor-bound-complete}): the
reduction closes, closure and minimality are consistent, and Difference anchors, the
closure bounds, minimality completes.
\end{enumerate}

The foundation of The Actualization Theory is therefore complete: the primitive Difference,
anchored by the Difference Principle; the closure that bounds its dynamics at $N=96$; and
the minimal theory that compresses the entire derivation structure to the irreducible
hierarchy Difference → Actualization → Spectrum → Physics. The following chapter derives,
from this foundation, the inevitable spectrum itself.

\begin{thebibliography}{99}
\bibitem{QG267} TQM-QG Phase 267, \emph{Conservation Principle Audit}, Docs/Research.
\bibitem{QG268} TQM-QG Phase 268, \emph{Count Conservation Origin}, Docs/Research.
\bibitem{QG278} TQM-QG Phase 278, \emph{Fundamental Boundary Audit}, Docs/Research.
\bibitem{QG282} TQM-QG Phase 282, \emph{Boundary Origin Audit} (Closure Principle), Docs/Research.
\bibitem{QG292} TQM-QG Phase 292, \emph{Foundation Stress Test}, Docs/Research.
\bibitem{QG293} TQM-QG Phase 293, \emph{Hierarchy Necessity Audit}, Docs/Research.
\bibitem{QG294} TQM-QG Phase 294, \emph{Minimal Theory Audit}, Docs/Research.
\bibitem{MONO006} TQM-MONO006, \emph{A01 Actualization Origin Audit}, Docs/Zenodo/Monograph.
\bibitem{MONO007} TQM-MONO007, \emph{Referee-Readiness Resolution}, Docs/Zenodo/Monograph.
\end{thebibliography}
 after the preamble
%         that defines the theorem environments (or include via the master file).
% ======================================================================

% ---- Theorem environments (define in the master preamble if not present) ----
% \newtheorem{definition}{Definition}[chapter]
% \newtheorem{lemma}[definition]{Lemma}
% \newtheorem{theorem}[definition]{Theorem}
% \newtheorem{corollary}[definition]{Corollary}
% \newtheorem{remark}[definition]{Remark}

\chapter{Closure and the Minimal Theory}
\label{chap:closure-minimal}

\section{Introduction}
\label{sec:closure-introduction}

In the preceding chapters we established the primitives $\{\text{Difference},\eta\}$
(Chapters~\ref{chap:difference} and \ref{chap:eta}) and the derived process face of
Difference, Actualization, whose convergence yields the inevitable spectrum
(Chapter~\ref{chap:actualization}). This chapter completes the foundational picture by
establishing two structural results. First, the \emph{Closure Principle}: the boundary of the
theory is the fixed point of Difference's own dynamics, and --- decisively --- the Closure
Principle \emph{characterizes} the fundamental boundary but does \emph{not derive} the
primitive \cite{QG278,QG282}. Second, the \emph{minimal theory}: the hierarchy
\begin{equation}
\label{eq:minimal}
\text{Difference}\;\longrightarrow\;\text{Actualization}\;\longrightarrow\;\text{Inevitable Spectrum}\;\longrightarrow\;\text{Physics}
\end{equation}
is complete and irreducible \cite{QG293,QG294}: every intermediate layer is derivable, and
no further primitive is required.

The results of this chapter are the canonical end-state of the closure and reduction program
\cite{QG267,QG268,QG278,QG282,QG293,QG294}, consistent with the MONO006 resolution and the
MONO007 referee-ready architecture \cite{MONO006,MONO007}. Throughout we use only accepted
canonical results and introduce no new physics and no speculation.

\section{The Closure Principle}
\label{sec:closure-principle}

\begin{definition}[Closure Principle]
\label{def:closure-principle}
The \emph{Closure Principle} is the statement that the boundary of the theory is the fixed
point of the actualization process: the primitive is the process, and the boundary is its
converged fixed point \cite{QG282}.
\end{definition}

\begin{theorem}[The Closure Principle holds]
\label{thm:closure-holds}
The boundary of the theory is the fixed point of Difference's own count-producing dynamics:
the topology saturates at zero residual change, and every initial pattern converges to the
same final geometry.
\end{theorem}

\begin{proof}
The boundary-origin audit \cite{QG282} establishes the closure: the actualization dynamics
has a self-reinforcing feedback bounded by network capacity, so positive feedback saturates
at a stable fixed point with zero residual link growth and a content-independent final
geometry. The observable sector is therefore the converged attractor of the dynamics, and
the boundary is the closure of that dynamics. The closure is derived --- it is an output of
the process --- not an imported cut-off.
\end{proof}

\begin{corollary}[The closure is \texorpdfstring{$N=96$}{N=96}]
\label{cor:closure-n96}
The fixed point of the actualization process is the $N=96$ network, whose eigenspectrum is
the inevitable D96 spectrum.
\end{corollary}

\begin{proof}
By Theorem~\ref{thm:closure-holds}, the boundary is the converged attractor of the
dynamics; by Theorem~\ref{thm:n96} of Chapter~\ref{chap:actualization}, that attractor is
the $N=96$ network; and by Corollary~\ref{cor:spectrum-eigenspectrum} of
Chapter~\ref{chap:actualization}, its eigenspectrum is the inevitable D96 spectrum.
\end{proof}

\section{Boundary and Fixed Point}
\label{sec:boundary-fixed-point}

\begin{definition}[Boundary]
\label{def:boundary}
A \emph{boundary} of the theory is the configuration at which the derived structure closes:
the point where the count-producing dynamics has converged and no further change occurs.
\end{definition}

\begin{theorem}[The boundary is the fixed point]
\label{thm:boundary-is-fixed}
The boundary of the theory is the stable fixed point of the actualization process; it is
derived, not primitive.
\end{theorem}

\begin{proof}
By Definition~\ref{def:closure-principle} and Theorem~\ref{thm:closure-holds}, the boundary
is the converged fixed point of the process. The derivation structure of the process is
established by Theorem~\ref{thm:derived-not-primitive} of
Chapter~\ref{chap:actualization}: the process is the face of Difference, and its fixed
point is an output of the dynamics. Hence the boundary is derived --- it is not an imported
input or a separate primitive.
\end{proof}

\begin{lemma}[The boundary characterizes the primitive]
\label{lem:boundary-characterizes}
The fixed-point boundary is a characterization of the primitive Difference: it specifies
how Difference's dynamics closes, and thereby \emph{what} the boundary is, without defining
\emph{what Difference is}.
\end{lemma}

\begin{proof}
The boundary is the fixed point of the process that Difference generates
(Theorem~\ref{thm:boundary-is-fixed}). It tells us where the process converges --- the
closure $N=96$ --- but presupposes the process and hence the primitive. A characterization
specifies properties of an object already given; it does not supply the object itself. The
fixed-point boundary therefore characterizes Difference without deriving it.
\end{proof}

\begin{corollary}[Closure does not derive Difference]
\label{cor:closure-not-derive}
The Closure Principle does not derive the primitive Difference. It characterizes the
boundary of the theory as the fixed point of Difference's own dynamics.
\end{corollary}

\begin{proof}
By Lemma~\ref{lem:boundary-characterizes}, the closure is a characterization, not a
derivation. Moreover, by Theorem~\ref{thm:fundamental-boundary} of
Chapter~\ref{chap:difference}, Difference is the fundamental boundary: irreducible, with its
count-conservation law as its definitional identity. A derivation of the primitive would
presuppose a deeper notion, which the fundamental-boundary audit \cite{QG278} shows does
not exist --- every attempt to reduce Difference reintroduces it. Hence the Closure
Principle characterizes, and does not derive, the primitive.
\end{proof}

\begin{remark}[Referee-safe reading]
Consistent with MONO007 \cite{MONO007}, every statement in this chapter that invokes the
Closure Principle is worded so that the principle is never presented as a derivation of
Difference: the primitive is given, and the closure characterizes where its process closes.
This removes the circularity objection A02 of the hostile-referee audit.
\end{remark}

\section{Self-Consistency}
\label{sec:self-consistency}

\begin{definition}[Self-consistency]
\label{def:self-consistency}
\emph{Self-consistency} is the requirement that the theory be free of internal
contradiction: non-contradiction requires comparing parts --- a difference.
\end{definition}

\begin{theorem}[Self-consistency presupposes Difference]
\label{thm:self-consistency-presupposes}
Self-consistency is derivable from Difference: non-contradiction is the comparison of parts,
and identity is the zero difference.
\end{theorem}

\begin{proof}
The fundamental-boundary audit \cite{QG278} establishes that self-consistency presupposes
Difference: non-contradiction requires comparing parts of the theory --- a difference ---
and identity is the zero difference \cite{QG268}. The conservation of the count, which is
the definitional identity of the primitive, is the universal conservation principle of which
all conservation laws are manifestations \cite{QG267}. Hence self-consistency is not an
independent postulate but a consequence of the primitive's structure.
\end{proof}

\begin{corollary}[Conservation is universal]
\label{cor:conservation-universal}
All conservation laws of the theory --- norm, trace, count, Bianchi, Noether currents ---
are manifestations of one universal conservation principle rooted in the identity of
Difference.
\end{corollary}

\begin{proof}
The conservation-principle audit \cite{QG267} verifies that the six conservation laws (Born
norm, unitarity, trace conservation, count conservation, Bianchi conservation, Noether
currents) are all manifestations of one deeper principle. By
Theorem~\ref{thm:count-conservation} of Chapter~\ref{chap:actualization}, count
conservation is the definitional identity of the primitive. Hence the universal conservation
principle is the identity of Difference, and all conservation laws are its manifestations.
\end{proof}

\section{Individuation}
\label{sec:individuation}

\begin{definition}[Individuation]
\label{def:individuation}
\emph{Individuation} is the process by which the one Difference gives rise to distinct
sectors: the differentiation of the theory's structure into the separate domains of
observable physics.
\end{definition}

\begin{theorem}[Individuation derives from Difference]
\label{thm:individuation}
The individuation of the theory into distinct sectors is derivable from Difference: distinct
sectors are distinct differences.
\end{theorem}

\begin{proof}
Sectors are distinguished by their different spectral access --- the projection classes over
the operator layer established in Chapter~\ref{chap:actualization}
(Theorem~\ref{thm:resonance-readout}). Each sector is a distinct pattern of difference in
the count structure. The individuation of the theory is therefore the differentiation of the
one Difference into its distinct faces and readouts; it is not an independent primitive but
a derived structure of the single difference principle.
\end{proof}

\section{The Difference Principle}
\label{sec:difference-principle}

\begin{definition}[Difference Principle]
\label{def:difference-principle}
The \emph{Difference Principle} is the statement that Difference is the root of the
hierarchy: the first concept, presupposed by everything else, to which every reduction
attempt returns.
\end{definition}

\begin{theorem}[Difference is the first concept]
\label{thm:difference-first}
Difference is the first concept of the theory: every attempt to reduce it reintroduces it,
so it is a self-referential fundamental boundary.
\end{theorem}

\begin{proof}
The fundamental-boundary audit \cite{QG278} performs the self-referential check. Actualization
presupposes Difference (an act of actualizing is a change --- a before/after difference; a
Q-event is a transition). Question presupposes Difference (a question is a gap --- the
difference between known and unknown). Self-consistency presupposes Difference
(non-contradiction requires comparing parts; identity is the zero difference). And
Difference itself is the first concept: any attempt to reduce it reintroduces it --- "two
distinct things" uses distinct = difference, "a boundary" is where things differ, "a
transition" is a before/after difference, "empty vs non-empty" is itself a difference. Hence
Difference is a self-referential fundamental boundary, and the Difference Principle states
exactly this: Difference is the root.
\end{proof}

\begin{corollary}[The Difference Principle anchors the hierarchy]
\label{cor:difference-anchors}
The Difference Principle anchors the canonical hierarchy
\eqref{eq:minimal}: every layer --- Actualization, Spectrum, Physics --- presupposes
Difference as its root.
\end{corollary}

\begin{proof}
By Theorem~\ref{thm:difference-first}, every reduction returns to Difference; by
Theorem~\ref{thm:process-face} of Chapter~\ref{chap:actualization}, Actualization is the
process face of Difference; and the spectrum and physics are the outputs of that process
(Chapter~\ref{chap:actualization}, Theorem~\ref{thm:spectrum-inevitable}). Hence the
Difference Principle anchors the entire hierarchy at its root.
\end{proof}

\section{The Minimal Theory}
\label{sec:minimal-theory}

\begin{theorem}[The minimal hierarchy]
\label{thm:minimal-hierarchy}
The minimal hierarchy is
\[
\text{Difference}\;\longrightarrow\;\text{Actualization}\;\longrightarrow\;\text{Inevitable Spectrum}\;\longrightarrow\;\text{Physics}.
\]
It is complete: every phase of the theory can be rewritten using only this hierarchy.
\end{theorem}

\begin{proof}
The hierarchy-necessity audit \cite{QG293} removes each intermediate layer of the fuller
hierarchy (Difference → Actualization → Closure → Conservation → Resonance → Spectrum →
Question → Measurement → Physics) and classifies it. Actualization is INDISPENSABLE (the
count-producing process). Closure is COMPRESSIBLE into Actualization (the fixed point of the
process). Conservation is COMPRESSIBLE into Difference/network (the definitional identity
plus the graph handshake). Resonance is COMPRESSIBLE into Spectrum (the operator readout).
The minimal-theory audit \cite{QG294} then re-derives every compressed layer from the
minimal hierarchy, classifying each DERIVABLE or ACTUALLY REQUIRED, and finds
\emph{5 DERIVABLE / 0 ACTUALLY REQUIRED}. Hence the minimal hierarchy is complete.
\end{proof}

\begin{theorem}[The minimal theory is irreducible]
\label{thm:minimal-irreducible}
The minimal theory is irreducible: no layer of the hierarchy
\eqref{eq:minimal} can be removed without losing derived content, and no further primitive
exists.
\end{theorem}

\begin{proof}
Irreducibility has two parts. \emph{Primitives:} by Theorem~\ref{thm:minimal-foundation} of
Chapter~\ref{chap:eta}, the minimal primitive set is $\{\text{Difference},\eta\}$; by
Theorem~\ref{thm:difference-first}, no deeper notion exists; by
Theorem~\ref{thm:derived-not-primitive} of Chapter~\ref{chap:actualization}, Actualization
is derived, not primitive. \emph{Layers:} the foundation stress test \cite{QG292}
establishes that removing Difference collapses all layers, and the hierarchy-necessity audit
\cite{QG293} establishes that removing Actualization destroys the count-producing process.
The minimal-theory audit \cite{QG294} confirms that the four layers of \eqref{eq:minimal}
are exactly what is required, with zero additional inputs. Hence the minimal theory is
irreducible.
\end{proof}

\begin{corollary}[The canonical architecture]
\label{cor:minimal-architecture}
The canonical architecture of the theory is the minimal hierarchy
\eqref{eq:minimal}: Difference → Actualization → Inevitable Spectrum → Physics, with
primitives $\{\text{Difference},\eta\}$, Actualization derived, the spectrum inevitable, and
physics emergent.
\end{corollary}

\begin{proof}
The minimal hierarchy is complete (Theorem~\ref{thm:minimal-hierarchy}) and irreducible
(Theorem~\ref{thm:minimal-irreducible}); the primitives are the pair
$\{\text{Difference},\eta\}$ (Theorem~\ref{thm:minimal-foundation} of
Chapter~\ref{chap:eta}); Actualization is derived (Theorem~\ref{thm:derived-not-primitive}
of Chapter~\ref{chap:actualization}); the spectrum is inevitable
(Theorem~\ref{thm:spectrum-inevitable} of Chapter~\ref{chap:actualization}). This is the
referee-ready structure of MONO007 \cite{MONO007}.
\end{proof}

\begin{remark}[Referee-safe completeness]
Consistent with MONO007 \cite{MONO007}, where the monograph claims completeness of the
minimal theory, it is understood in the qualified sense: every \emph{observable} is derived
from the minimal hierarchy, while the standard-model \emph{Lagrangian} and its dynamics are
\emph{hosted} --- a boundary, not a derivation gap. The minimal theory is irreducible as a
derivation structure; it does not assert a derivation of every aspect of the hosted dynamics.
\end{remark}

\section{Consequences}
\label{sec:closure-consequences}

\begin{lemma}[The reduction closes at the minimal hierarchy]
\label{lem:reduction-closes}
The reduction program closes at the minimal hierarchy \eqref{eq:minimal}: all intermediate
layers of the fuller hierarchy are derivable from it, and none is actually required as an
input.
\end{lemma}

\begin{proof}
The minimal-theory audit \cite{QG294} classifies the five compressed layers --- Closure,
Conservation, Resonance, and the remaining intermediates --- as DERIVABLE, with zero
ACTUALLY REQUIRED. The re-derivation uses only the minimal hierarchy. Hence the reduction
closes at \eqref{eq:minimal}.
\end{proof}

\begin{theorem}[Closure and minimality are consistent]
\label{thm:closure-minimality}
The Closure Principle and the minimal theory are mutually consistent: the closure fixed
point $N=96$ is the output of the minimal hierarchy, and the minimal hierarchy requires the
closure to generate the spectrum.
\end{theorem}

\begin{proof}
The closure fixed point $N=96$ is the output of the actualization process
(Corollary~\ref{cor:closure-n96}), which is a layer of the minimal hierarchy. Conversely,
the spectrum layer of the minimal hierarchy is the eigenspectrum of the converged network
(Chapter~\ref{chap:actualization}, Corollary~\ref{cor:spectrum-eigenspectrum}), so the
minimal hierarchy requires the closure to produce the spectrum. The two results are the
same structure viewed from the process side (closure) and the output side (spectrum).
\end{proof}

\begin{corollary}[Difference anchors; closure bounds; minimality completes]
\label{cor:anchor-bound-complete}
In the canonical architecture, Difference anchors the hierarchy (the root), the closure
bounds it (the fixed point), and the minimal theory completes it (the irreducible
derivation structure).
\end{corollary}

\begin{proof}
Difference anchors by the Difference Principle (Theorem~\ref{thm:difference-first},
Corollary~\ref{cor:difference-anchors}); the closure bounds by the Closure Principle
(Theorem~\ref{thm:closure-holds}, Corollary~\ref{cor:closure-n96}); the minimal theory
completes by Theorem~\ref{thm:minimal-hierarchy} and Theorem~\ref{thm:minimal-irreducible}.
The three roles are distinct and jointly exhaustive of the foundational structure.
\end{proof}

\section{Summary}
\label{sec:closure-summary}

This chapter has established the foundational closure of the theory. We have proved:

\begin{enumerate}
\item \textbf{The Closure Principle} (Theorem~\ref{thm:closure-holds},
Corollary~\ref{cor:closure-n96}): the boundary of the theory is the fixed point of
Difference's own dynamics, closing at $N=96$;
\item \textbf{Boundary and fixed point} (Theorem~\ref{thm:boundary-is-fixed},
Lemma~\ref{lem:boundary-characterizes}, Corollary~\ref{cor:closure-not-derive}): the
boundary is derived, characterizes the primitive, and --- decisively --- the Closure
Principle does \emph{not} derive Difference;
\item \textbf{Self-consistency} (Theorem~\ref{thm:self-consistency-presupposes},
Corollary~\ref{cor:conservation-universal}): self-consistency presupposes Difference, and
conservation is universal;
\item \textbf{Individuation} (Theorem~\ref{thm:individuation}): the distinct sectors derive
from Difference;
\item \textbf{The Difference Principle} (Theorem~\ref{thm:difference-first},
Corollary~\ref{cor:difference-anchors}): Difference is the self-referential first concept,
anchoring the hierarchy;
\item \textbf{The minimal theory} (Theorem~\ref{thm:minimal-hierarchy},
Theorem~\ref{thm:minimal-irreducible}, Corollary~\ref{cor:minimal-architecture}): the
minimal hierarchy is complete and irreducible;
\item \textbf{Consequences} (Lemma~\ref{lem:reduction-closes},
Theorem~\ref{thm:closure-minimality}, Corollary~\ref{cor:anchor-bound-complete}): the
reduction closes, closure and minimality are consistent, and Difference anchors, the
closure bounds, minimality completes.
\end{enumerate}

The foundation of The Actualization Theory is therefore complete: the primitive Difference,
anchored by the Difference Principle; the closure that bounds its dynamics at $N=96$; and
the minimal theory that compresses the entire derivation structure to the irreducible
hierarchy Difference → Actualization → Spectrum → Physics. The following chapter derives,
from this foundation, the inevitable spectrum itself.

\begin{thebibliography}{99}
\bibitem{QG267} TQM-QG Phase 267, \emph{Conservation Principle Audit}, Docs/Research.
\bibitem{QG268} TQM-QG Phase 268, \emph{Count Conservation Origin}, Docs/Research.
\bibitem{QG278} TQM-QG Phase 278, \emph{Fundamental Boundary Audit}, Docs/Research.
\bibitem{QG282} TQM-QG Phase 282, \emph{Boundary Origin Audit} (Closure Principle), Docs/Research.
\bibitem{QG292} TQM-QG Phase 292, \emph{Foundation Stress Test}, Docs/Research.
\bibitem{QG293} TQM-QG Phase 293, \emph{Hierarchy Necessity Audit}, Docs/Research.
\bibitem{QG294} TQM-QG Phase 294, \emph{Minimal Theory Audit}, Docs/Research.
\bibitem{MONO006} TQM-MONO006, \emph{A01 Actualization Origin Audit}, Docs/Zenodo/Monograph.
\bibitem{MONO007} TQM-MONO007, \emph{Referee-Readiness Resolution}, Docs/Zenodo/Monograph.
\end{thebibliography}
 after the preamble
%         that defines the theorem environments (or include via the master file).
% ======================================================================

% ---- Theorem environments (define in the master preamble if not present) ----
% \newtheorem{definition}{Definition}[chapter]
% \newtheorem{lemma}[definition]{Lemma}
% \newtheorem{theorem}[definition]{Theorem}
% \newtheorem{corollary}[definition]{Corollary}
% \newtheorem{remark}[definition]{Remark}

\chapter{Closure and the Minimal Theory}
\label{chap:closure-minimal}

\section{Introduction}
\label{sec:closure-introduction}

In the preceding chapters we established the primitives $\{\text{Difference},\eta\}$
(Chapters~\ref{chap:difference} and \ref{chap:eta}) and the derived process face of
Difference, Actualization, whose convergence yields the inevitable spectrum
(Chapter~\ref{chap:actualization}). This chapter completes the foundational picture by
establishing two structural results. First, the \emph{Closure Principle}: the boundary of the
theory is the fixed point of Difference's own dynamics, and --- decisively --- the Closure
Principle \emph{characterizes} the fundamental boundary but does \emph{not derive} the
primitive \cite{QG278,QG282}. Second, the \emph{minimal theory}: the hierarchy
\begin{equation}
\label{eq:minimal}
\text{Difference}\;\longrightarrow\;\text{Actualization}\;\longrightarrow\;\text{Inevitable Spectrum}\;\longrightarrow\;\text{Physics}
\end{equation}
is complete and irreducible \cite{QG293,QG294}: every intermediate layer is derivable, and
no further primitive is required.

The results of this chapter are the canonical end-state of the closure and reduction program
\cite{QG267,QG268,QG278,QG282,QG293,QG294}, consistent with the MONO006 resolution and the
MONO007 referee-ready architecture \cite{MONO006,MONO007}. Throughout we use only accepted
canonical results and introduce no new physics and no speculation.

\section{The Closure Principle}
\label{sec:closure-principle}

\begin{definition}[Closure Principle]
\label{def:closure-principle}
The \emph{Closure Principle} is the statement that the boundary of the theory is the fixed
point of the actualization process: the primitive is the process, and the boundary is its
converged fixed point \cite{QG282}.
\end{definition}

\begin{theorem}[The Closure Principle holds]
\label{thm:closure-holds}
The boundary of the theory is the fixed point of Difference's own count-producing dynamics:
the topology saturates at zero residual change, and every initial pattern converges to the
same final geometry.
\end{theorem}

\begin{proof}
The boundary-origin audit \cite{QG282} establishes the closure: the actualization dynamics
has a self-reinforcing feedback bounded by network capacity, so positive feedback saturates
at a stable fixed point with zero residual link growth and a content-independent final
geometry. The observable sector is therefore the converged attractor of the dynamics, and
the boundary is the closure of that dynamics. The closure is derived --- it is an output of
the process --- not an imported cut-off.
\end{proof}

\begin{corollary}[The closure is \texorpdfstring{$N=96$}{N=96}]
\label{cor:closure-n96}
The fixed point of the actualization process is the $N=96$ network, whose eigenspectrum is
the inevitable D96 spectrum.
\end{corollary}

\begin{proof}
By Theorem~\ref{thm:closure-holds}, the boundary is the converged attractor of the
dynamics; by Theorem~\ref{thm:n96} of Chapter~\ref{chap:actualization}, that attractor is
the $N=96$ network; and by Corollary~\ref{cor:spectrum-eigenspectrum} of
Chapter~\ref{chap:actualization}, its eigenspectrum is the inevitable D96 spectrum.
\end{proof}

\section{Boundary and Fixed Point}
\label{sec:boundary-fixed-point}

\begin{definition}[Boundary]
\label{def:boundary}
A \emph{boundary} of the theory is the configuration at which the derived structure closes:
the point where the count-producing dynamics has converged and no further change occurs.
\end{definition}

\begin{theorem}[The boundary is the fixed point]
\label{thm:boundary-is-fixed}
The boundary of the theory is the stable fixed point of the actualization process; it is
derived, not primitive.
\end{theorem}

\begin{proof}
By Definition~\ref{def:closure-principle} and Theorem~\ref{thm:closure-holds}, the boundary
is the converged fixed point of the process. The derivation structure of the process is
established by Theorem~\ref{thm:derived-not-primitive} of
Chapter~\ref{chap:actualization}: the process is the face of Difference, and its fixed
point is an output of the dynamics. Hence the boundary is derived --- it is not an imported
input or a separate primitive.
\end{proof}

\begin{lemma}[The boundary characterizes the primitive]
\label{lem:boundary-characterizes}
The fixed-point boundary is a characterization of the primitive Difference: it specifies
how Difference's dynamics closes, and thereby \emph{what} the boundary is, without defining
\emph{what Difference is}.
\end{lemma}

\begin{proof}
The boundary is the fixed point of the process that Difference generates
(Theorem~\ref{thm:boundary-is-fixed}). It tells us where the process converges --- the
closure $N=96$ --- but presupposes the process and hence the primitive. A characterization
specifies properties of an object already given; it does not supply the object itself. The
fixed-point boundary therefore characterizes Difference without deriving it.
\end{proof}

\begin{corollary}[Closure does not derive Difference]
\label{cor:closure-not-derive}
The Closure Principle does not derive the primitive Difference. It characterizes the
boundary of the theory as the fixed point of Difference's own dynamics.
\end{corollary}

\begin{proof}
By Lemma~\ref{lem:boundary-characterizes}, the closure is a characterization, not a
derivation. Moreover, by Theorem~\ref{thm:fundamental-boundary} of
Chapter~\ref{chap:difference}, Difference is the fundamental boundary: irreducible, with its
count-conservation law as its definitional identity. A derivation of the primitive would
presuppose a deeper notion, which the fundamental-boundary audit \cite{QG278} shows does
not exist --- every attempt to reduce Difference reintroduces it. Hence the Closure
Principle characterizes, and does not derive, the primitive.
\end{proof}

\begin{remark}[Referee-safe reading]
Consistent with MONO007 \cite{MONO007}, every statement in this chapter that invokes the
Closure Principle is worded so that the principle is never presented as a derivation of
Difference: the primitive is given, and the closure characterizes where its process closes.
This removes the circularity objection A02 of the hostile-referee audit.
\end{remark}

\section{Self-Consistency}
\label{sec:self-consistency}

\begin{definition}[Self-consistency]
\label{def:self-consistency}
\emph{Self-consistency} is the requirement that the theory be free of internal
contradiction: non-contradiction requires comparing parts --- a difference.
\end{definition}

\begin{theorem}[Self-consistency presupposes Difference]
\label{thm:self-consistency-presupposes}
Self-consistency is derivable from Difference: non-contradiction is the comparison of parts,
and identity is the zero difference.
\end{theorem}

\begin{proof}
The fundamental-boundary audit \cite{QG278} establishes that self-consistency presupposes
Difference: non-contradiction requires comparing parts of the theory --- a difference ---
and identity is the zero difference \cite{QG268}. The conservation of the count, which is
the definitional identity of the primitive, is the universal conservation principle of which
all conservation laws are manifestations \cite{QG267}. Hence self-consistency is not an
independent postulate but a consequence of the primitive's structure.
\end{proof}

\begin{corollary}[Conservation is universal]
\label{cor:conservation-universal}
All conservation laws of the theory --- norm, trace, count, Bianchi, Noether currents ---
are manifestations of one universal conservation principle rooted in the identity of
Difference.
\end{corollary}

\begin{proof}
The conservation-principle audit \cite{QG267} verifies that the six conservation laws (Born
norm, unitarity, trace conservation, count conservation, Bianchi conservation, Noether
currents) are all manifestations of one deeper principle. By
Theorem~\ref{thm:count-conservation} of Chapter~\ref{chap:actualization}, count
conservation is the definitional identity of the primitive. Hence the universal conservation
principle is the identity of Difference, and all conservation laws are its manifestations.
\end{proof}

\section{Individuation}
\label{sec:individuation}

\begin{definition}[Individuation]
\label{def:individuation}
\emph{Individuation} is the process by which the one Difference gives rise to distinct
sectors: the differentiation of the theory's structure into the separate domains of
observable physics.
\end{definition}

\begin{theorem}[Individuation derives from Difference]
\label{thm:individuation}
The individuation of the theory into distinct sectors is derivable from Difference: distinct
sectors are distinct differences.
\end{theorem}

\begin{proof}
Sectors are distinguished by their different spectral access --- the projection classes over
the operator layer established in Chapter~\ref{chap:actualization}
(Theorem~\ref{thm:resonance-readout}). Each sector is a distinct pattern of difference in
the count structure. The individuation of the theory is therefore the differentiation of the
one Difference into its distinct faces and readouts; it is not an independent primitive but
a derived structure of the single difference principle.
\end{proof}

\section{The Difference Principle}
\label{sec:difference-principle}

\begin{definition}[Difference Principle]
\label{def:difference-principle}
The \emph{Difference Principle} is the statement that Difference is the root of the
hierarchy: the first concept, presupposed by everything else, to which every reduction
attempt returns.
\end{definition}

\begin{theorem}[Difference is the first concept]
\label{thm:difference-first}
Difference is the first concept of the theory: every attempt to reduce it reintroduces it,
so it is a self-referential fundamental boundary.
\end{theorem}

\begin{proof}
The fundamental-boundary audit \cite{QG278} performs the self-referential check. Actualization
presupposes Difference (an act of actualizing is a change --- a before/after difference; a
Q-event is a transition). Question presupposes Difference (a question is a gap --- the
difference between known and unknown). Self-consistency presupposes Difference
(non-contradiction requires comparing parts; identity is the zero difference). And
Difference itself is the first concept: any attempt to reduce it reintroduces it --- "two
distinct things" uses distinct = difference, "a boundary" is where things differ, "a
transition" is a before/after difference, "empty vs non-empty" is itself a difference. Hence
Difference is a self-referential fundamental boundary, and the Difference Principle states
exactly this: Difference is the root.
\end{proof}

\begin{corollary}[The Difference Principle anchors the hierarchy]
\label{cor:difference-anchors}
The Difference Principle anchors the canonical hierarchy
\eqref{eq:minimal}: every layer --- Actualization, Spectrum, Physics --- presupposes
Difference as its root.
\end{corollary}

\begin{proof}
By Theorem~\ref{thm:difference-first}, every reduction returns to Difference; by
Theorem~\ref{thm:process-face} of Chapter~\ref{chap:actualization}, Actualization is the
process face of Difference; and the spectrum and physics are the outputs of that process
(Chapter~\ref{chap:actualization}, Theorem~\ref{thm:spectrum-inevitable}). Hence the
Difference Principle anchors the entire hierarchy at its root.
\end{proof}

\section{The Minimal Theory}
\label{sec:minimal-theory}

\begin{theorem}[The minimal hierarchy]
\label{thm:minimal-hierarchy}
The minimal hierarchy is
\[
\text{Difference}\;\longrightarrow\;\text{Actualization}\;\longrightarrow\;\text{Inevitable Spectrum}\;\longrightarrow\;\text{Physics}.
\]
It is complete: every phase of the theory can be rewritten using only this hierarchy.
\end{theorem}

\begin{proof}
The hierarchy-necessity audit \cite{QG293} removes each intermediate layer of the fuller
hierarchy (Difference → Actualization → Closure → Conservation → Resonance → Spectrum →
Question → Measurement → Physics) and classifies it. Actualization is INDISPENSABLE (the
count-producing process). Closure is COMPRESSIBLE into Actualization (the fixed point of the
process). Conservation is COMPRESSIBLE into Difference/network (the definitional identity
plus the graph handshake). Resonance is COMPRESSIBLE into Spectrum (the operator readout).
The minimal-theory audit \cite{QG294} then re-derives every compressed layer from the
minimal hierarchy, classifying each DERIVABLE or ACTUALLY REQUIRED, and finds
\emph{5 DERIVABLE / 0 ACTUALLY REQUIRED}. Hence the minimal hierarchy is complete.
\end{proof}

\begin{theorem}[The minimal theory is irreducible]
\label{thm:minimal-irreducible}
The minimal theory is irreducible: no layer of the hierarchy
\eqref{eq:minimal} can be removed without losing derived content, and no further primitive
exists.
\end{theorem}

\begin{proof}
Irreducibility has two parts. \emph{Primitives:} by Theorem~\ref{thm:minimal-foundation} of
Chapter~\ref{chap:eta}, the minimal primitive set is $\{\text{Difference},\eta\}$; by
Theorem~\ref{thm:difference-first}, no deeper notion exists; by
Theorem~\ref{thm:derived-not-primitive} of Chapter~\ref{chap:actualization}, Actualization
is derived, not primitive. \emph{Layers:} the foundation stress test \cite{QG292}
establishes that removing Difference collapses all layers, and the hierarchy-necessity audit
\cite{QG293} establishes that removing Actualization destroys the count-producing process.
The minimal-theory audit \cite{QG294} confirms that the four layers of \eqref{eq:minimal}
are exactly what is required, with zero additional inputs. Hence the minimal theory is
irreducible.
\end{proof}

\begin{corollary}[The canonical architecture]
\label{cor:minimal-architecture}
The canonical architecture of the theory is the minimal hierarchy
\eqref{eq:minimal}: Difference → Actualization → Inevitable Spectrum → Physics, with
primitives $\{\text{Difference},\eta\}$, Actualization derived, the spectrum inevitable, and
physics emergent.
\end{corollary}

\begin{proof}
The minimal hierarchy is complete (Theorem~\ref{thm:minimal-hierarchy}) and irreducible
(Theorem~\ref{thm:minimal-irreducible}); the primitives are the pair
$\{\text{Difference},\eta\}$ (Theorem~\ref{thm:minimal-foundation} of
Chapter~\ref{chap:eta}); Actualization is derived (Theorem~\ref{thm:derived-not-primitive}
of Chapter~\ref{chap:actualization}); the spectrum is inevitable
(Theorem~\ref{thm:spectrum-inevitable} of Chapter~\ref{chap:actualization}). This is the
referee-ready structure of MONO007 \cite{MONO007}.
\end{proof}

\begin{remark}[Referee-safe completeness]
Consistent with MONO007 \cite{MONO007}, where the monograph claims completeness of the
minimal theory, it is understood in the qualified sense: every \emph{observable} is derived
from the minimal hierarchy, while the standard-model \emph{Lagrangian} and its dynamics are
\emph{hosted} --- a boundary, not a derivation gap. The minimal theory is irreducible as a
derivation structure; it does not assert a derivation of every aspect of the hosted dynamics.
\end{remark}

\section{Consequences}
\label{sec:closure-consequences}

\begin{lemma}[The reduction closes at the minimal hierarchy]
\label{lem:reduction-closes}
The reduction program closes at the minimal hierarchy \eqref{eq:minimal}: all intermediate
layers of the fuller hierarchy are derivable from it, and none is actually required as an
input.
\end{lemma}

\begin{proof}
The minimal-theory audit \cite{QG294} classifies the five compressed layers --- Closure,
Conservation, Resonance, and the remaining intermediates --- as DERIVABLE, with zero
ACTUALLY REQUIRED. The re-derivation uses only the minimal hierarchy. Hence the reduction
closes at \eqref{eq:minimal}.
\end{proof}

\begin{theorem}[Closure and minimality are consistent]
\label{thm:closure-minimality}
The Closure Principle and the minimal theory are mutually consistent: the closure fixed
point $N=96$ is the output of the minimal hierarchy, and the minimal hierarchy requires the
closure to generate the spectrum.
\end{theorem}

\begin{proof}
The closure fixed point $N=96$ is the output of the actualization process
(Corollary~\ref{cor:closure-n96}), which is a layer of the minimal hierarchy. Conversely,
the spectrum layer of the minimal hierarchy is the eigenspectrum of the converged network
(Chapter~\ref{chap:actualization}, Corollary~\ref{cor:spectrum-eigenspectrum}), so the
minimal hierarchy requires the closure to produce the spectrum. The two results are the
same structure viewed from the process side (closure) and the output side (spectrum).
\end{proof}

\begin{corollary}[Difference anchors; closure bounds; minimality completes]
\label{cor:anchor-bound-complete}
In the canonical architecture, Difference anchors the hierarchy (the root), the closure
bounds it (the fixed point), and the minimal theory completes it (the irreducible
derivation structure).
\end{corollary}

\begin{proof}
Difference anchors by the Difference Principle (Theorem~\ref{thm:difference-first},
Corollary~\ref{cor:difference-anchors}); the closure bounds by the Closure Principle
(Theorem~\ref{thm:closure-holds}, Corollary~\ref{cor:closure-n96}); the minimal theory
completes by Theorem~\ref{thm:minimal-hierarchy} and Theorem~\ref{thm:minimal-irreducible}.
The three roles are distinct and jointly exhaustive of the foundational structure.
\end{proof}

\section{Summary}
\label{sec:closure-summary}

This chapter has established the foundational closure of the theory. We have proved:

\begin{enumerate}
\item \textbf{The Closure Principle} (Theorem~\ref{thm:closure-holds},
Corollary~\ref{cor:closure-n96}): the boundary of the theory is the fixed point of
Difference's own dynamics, closing at $N=96$;
\item \textbf{Boundary and fixed point} (Theorem~\ref{thm:boundary-is-fixed},
Lemma~\ref{lem:boundary-characterizes}, Corollary~\ref{cor:closure-not-derive}): the
boundary is derived, characterizes the primitive, and --- decisively --- the Closure
Principle does \emph{not} derive Difference;
\item \textbf{Self-consistency} (Theorem~\ref{thm:self-consistency-presupposes},
Corollary~\ref{cor:conservation-universal}): self-consistency presupposes Difference, and
conservation is universal;
\item \textbf{Individuation} (Theorem~\ref{thm:individuation}): the distinct sectors derive
from Difference;
\item \textbf{The Difference Principle} (Theorem~\ref{thm:difference-first},
Corollary~\ref{cor:difference-anchors}): Difference is the self-referential first concept,
anchoring the hierarchy;
\item \textbf{The minimal theory} (Theorem~\ref{thm:minimal-hierarchy},
Theorem~\ref{thm:minimal-irreducible}, Corollary~\ref{cor:minimal-architecture}): the
minimal hierarchy is complete and irreducible;
\item \textbf{Consequences} (Lemma~\ref{lem:reduction-closes},
Theorem~\ref{thm:closure-minimality}, Corollary~\ref{cor:anchor-bound-complete}): the
reduction closes, closure and minimality are consistent, and Difference anchors, the
closure bounds, minimality completes.
\end{enumerate}

The foundation of The Actualization Theory is therefore complete: the primitive Difference,
anchored by the Difference Principle; the closure that bounds its dynamics at $N=96$; and
the minimal theory that compresses the entire derivation structure to the irreducible
hierarchy Difference → Actualization → Spectrum → Physics. The following chapter derives,
from this foundation, the inevitable spectrum itself.

\begin{thebibliography}{99}
\bibitem{QG267} TQM-QG Phase 267, \emph{Conservation Principle Audit}, Docs/Research.
\bibitem{QG268} TQM-QG Phase 268, \emph{Count Conservation Origin}, Docs/Research.
\bibitem{QG278} TQM-QG Phase 278, \emph{Fundamental Boundary Audit}, Docs/Research.
\bibitem{QG282} TQM-QG Phase 282, \emph{Boundary Origin Audit} (Closure Principle), Docs/Research.
\bibitem{QG292} TQM-QG Phase 292, \emph{Foundation Stress Test}, Docs/Research.
\bibitem{QG293} TQM-QG Phase 293, \emph{Hierarchy Necessity Audit}, Docs/Research.
\bibitem{QG294} TQM-QG Phase 294, \emph{Minimal Theory Audit}, Docs/Research.
\bibitem{MONO006} TQM-MONO006, \emph{A01 Actualization Origin Audit}, Docs/Zenodo/Monograph.
\bibitem{MONO007} TQM-MONO007, \emph{Referee-Readiness Resolution}, Docs/Zenodo/Monograph.
\end{thebibliography}
